\documentclass[11pt,a4paper]{article}
\usepackage[T1]{fontenc}
\usepackage[utf8]{inputenc}
\usepackage[french]{babel}
\usepackage{lmodern}
\usepackage{amsmath,amssymb}
\usepackage[margin=2.35cm]{geometry}
\usepackage{booktabs,longtable,tabularx}
\usepackage{enumitem}
\usepackage{xcolor}
\usepackage{microtype}
\usepackage{hyperref}

\hypersetup{colorlinks=true,linkcolor=blue!45!black,urlcolor=blue!45!black}
\setlist{nosep}
\newcommand{\code}[1]{\texttt{#1}}
\newcommand{\NW}{\mathrm{NW}}
\newcommand{\E}{\mathbb{E}}
\newcommand{\Kaut}{K^{*}_{\mathrm{aut}}}

\title{Spécification mathématique du modèle M4.2\\
\large société de crédit à production concave généralisée $F_\gamma(K)=A\,K^\gamma$}
\author{Programme M4.2 --- dossier \code{m4\_2\_credit\_soc}}
\date{27 juillet 2026}

\begin{document}
\maketitle

\begin{abstract}
M4.2 est le successeur direct de M4B : une société d'entités économiques
échangeant une ressource unique (le joule), avec capital productif réel $K$,
réseau de prêts nominaux perpétuels, chocs multiplicatifs individuels,
service d'intérêts, dépréciation et faillites en cascade par perte intégrale
des créances et destruction du capital résiduel. M4.2 généralise la
production de M4B, $F(K)=\sqrt{K}$, en $F_\gamma(K)=A\,K^\gamma$ avec
$0<\gamma<1$ et $A=1$ dans le modèle principal, dérive les taux d'intérêt
des rendements marginaux correspondants, fixe le pool d'appariement à deux
($k\equiv2$, retiré de la configuration) et nomme explicitement la fonction
d'intensité du marché $\eta$, dont la baseline est l'identité. Ce document
spécifie la mécanique exécutée par le moteur \code{m4\_2/model.py}
(version \code{m4\_2-1}), démontre les identités porteuses, documente les
tolérances numériques et établit la parité mesurée du cas $\gamma=1/2$ avec
M4B configuré à $k=2$. Le code exécuté fait foi ; ce document le décrit.
\end{abstract}

\tableofcontents

\section{Positionnement et notations}

M4.2 répond à la question de recherche du cahier des charges
(\code{prompts/PROMPT\_M4\_2.md}) : la pente $\tau$ de la distribution des
tailles d'avalanches de faillites peut-elle devenir une propriété modulable
par l'exposant de concavité $\gamma$, sans détruire le mécanisme endogène
d'accumulation--relaxation hérité de M4/M4B ?

Les symboles sont réservés ainsi, sans collision :
\begin{itemize}
\item $\gamma$ : exposant de concavité productive, $0<\gamma<1$ ;
\item $A$ : échelle productive, $A=1$ dans le modèle principal (elle ne
      varie que dans l'ablation de contrôle d'échelle, §\ref{sec:anorm}) ;
\item $\eta$ : fonction d'intensité du marché (jamais le choc) ;
\item $\xi_{i,t}$ : choc multiplicatif individuel ;
\item $\lambda$ : intensité des naissances de Poisson ;
\item $\delta$ : taux de dépréciation ; $\sigma$ : volatilité du choc ;
      $K_0$ : capital de naissance ;
\item $\tau$ : exposant statistique de la distribution des avalanches
      (jamais utilisé pour un temps) ;
\item $s_c$ : coupure exponentielle de taille finie ; $s_{\min}$ : seuil
      bas des ajustements de queue.
\end{itemize}
Le symbole $\alpha$ n'est employé ni pour l'échelle productive ni pour
l'exposant d'avalanche dans M4.2 (dans le code d'analyse hérité, le champ
\code{alpha} des ajustements désigne $\tau$ ou l'exposant tronqué ; le
rapport de campagne traduit systématiquement).

\section{État du modèle}

\subsection{Entités}

Une entité $i$ est identifiée par un entier jamais réutilisé. Son état
persistant est son capital productif réel $K_i \ge 0$ (en joules), son pas
de naissance, et son statut (vivante ou morte). S'y ajoutent des mesures de
pas remises à zéro en début de pas : production $\Pi_i$, intérêts reçus et
payés, et l'indicateur de défaut de service du pas courant.

\subsection{Contrats de prêt}

Le carnet (\code{LoanBook}) contient des contrats nominaux perpétuels
$(\ell, b, q, r)$ : la prêteuse $\ell$ détient une créance de principal $q$
sur l'emprunteuse $b$ au taux $r$ ; l'intérêt dû à chaque pas est $r\,q$,
sans échéance de remboursement du principal. Il existe au plus un contrat
par paire orientée $(\ell,b)$ : un nouveau prêt entre la même paire
\emph{fusionne} avec l'existant (principal additionné, taux moyen pondéré
par les principaux, §\ref{sec:merge}). Trois agrégats par entité sont
maintenus incrémentalement : créances $C_i=\sum_{b} q_{i\to b}$, dettes
$D_i=\sum_{\ell} q_{\ell\to i}$, et service dû
$S_i=\sum_{\ell} r_{\ell\to i}\,q_{\ell\to i}$.

La valeur nette est
\begin{equation}
\NW_i \;=\; K_i + C_i - D_i .
\label{eq:nw}
\end{equation}
Il n'existe ni dette d'amorçage, ni récupération, ni transfert de créances
lors d'une faillite.

\subsection{Configuration}

La configuration scientifique sérialisée (\code{Config}, dataclass gelée)
contient exactement : $\gamma$ (défaut $0{,}5$), $A$ (défaut $1{,}0$),
$\lambda$ (défaut $30$), $\delta$ (défaut $0{,}05$), $\sigma$ (défaut
$0{,}25$), $K_0$ (défaut $25$), la graine, l'horizon $T$ (défaut $2000$) et
le garde-fou de population \code{pop\_max} (défaut $30\,000$). Validation :
$0<\gamma<1$, $A>0$, $\lambda\ge0$, $0\le\delta<1$, $\sigma\ge0$, $K_0>0$,
$T\ge0$, $\code{pop\_max}\ge1$.

Les choix constitutifs, retirés de la configuration pour ne pas faire
passer des constantes pour des leviers scientifiques, sont des constantes
de module : le pool d'appariement $\code{POOL\_SIZE}=2$ ($k\equiv2$), les
tolérances numériques (§\ref{sec:tol}), et les règles institutionnelles
(moyenne géométrique, fusion, cancel+destroy, ordre des phases),
enregistrées dans \code{config.json} sous \code{fixed\_rules}.

\section{Chronologie exacte d'un pas}

L'ordre des huit phases est constitutif et identique à M4B :
naissances, choc, production, service des intérêts, dépréciation, marché du
crédit, faillites en cascade, mesures. Le générateur pseudo-aléatoire est
un \code{numpy.random.default\_rng(seed)} unique ; l'ordre de consommation
(Poisson des naissances, puis vecteur normal des chocs, puis tirages du
marché) est un invariant de reproductibilité.

\subsection{Phase 1 : naissances}

$B_t \sim \mathrm{Poisson}(\lambda)$ entités naissent avec $K=K_0$.
Avant les naissances, les mesures de pas de toutes les vivantes
(production, intérêts, indicateur de défaut) sont remises à zéro.
L'injection du pas vaut $B_t\,K_0$.

\subsection{Phase 2 : choc multiplicatif $\xi$}

Si $\sigma>0$, chaque vivante $i$ (dans l'ordre croissant des
identifiants) reçoit
\begin{equation}
\xi_{i,t}\sim\mathcal N\!\Bigl(-\tfrac{\sigma^2}{2},\,\sigma^2\Bigr),
\qquad K_i \leftarrow K_i\,e^{\xi_{i,t}},
\qquad \E\bigl[e^{\xi_{i,t}}\bigr]=1 .
\end{equation}
Le gain de choc agrégé (positif ou négatif) est mesuré pour le bilan.

\subsection{Phase 3 : production concave}
\label{sec:prod}

\begin{equation}
\Pi_i = F_\gamma(K_i) = A\,K_i^{\gamma},
\qquad K_i \leftarrow K_i + \Pi_i ,
\qquad 0<\gamma<1 .
\end{equation}
La production est entièrement retenue dans le capital. Le cas $\gamma=1/2$,
$A=1$ retrouve exactement la production $\sqrt{K}$ de M4B. À $K=0$,
$F_\gamma(0)=0$ sans garde particulière ($0^\gamma=0$ pour $\gamma>0$).

Le rendement marginal productif est
\begin{equation}
m_\gamma(K) = F_\gamma'(K) = A\,\gamma\,K^{\gamma-1},
\label{eq:marginal}
\end{equation}
décroissant en $K$ (concavité). Il n'est évalué que dans la phase de marché,
avec le plancher numérique de §\ref{sec:tol}.

\subsection{Phase 4 : service des intérêts}

Les emprunteuses sont parcourues dans l'ordre historique d'insertion de
leurs clés dans le carnet (invariant de reproductibilité hérité de M4B).
Pour l'emprunteuse $b$ de service dû $S_b$ et de capital disponible $K_b$ :
\begin{equation}
\varphi_b=\min\!\Bigl(1,\tfrac{K_b}{S_b}\Bigr),\qquad
K_b \leftarrow K_b - \varphi_b\,S_b ,
\end{equation}
et chaque prêteuse $\ell$ reçoit $\varphi_b\,r_{\ell b}\,q_{\ell b}$,
ajouté immédiatement à son capital (les paiements reçus avant son propre
tour de service sont donc disponibles pour payer). Si $\varphi_b<1$
(cas $S_b=0$ avec $K_b<0$ impossible ; le cas $S_b=0$ donne $\varphi_b=0$
uniquement si des contrats de service nul existaient, ce que la validation
du carnet exclut), l'entité est marquée en \emph{défaut de service} :
$K_b$ est mis à zéro, le paiement est réparti proportionnellement, et le
défaut la condamne en phase 7 même si sa valeur nette reste positive. Le
défaut n'annule ni le principal ni le service futur.

\subsection{Phase 5 : dépréciation}

$K_i \leftarrow (1-\delta)K_i$ pour toutes les vivantes ; un capital
inférieur à $10^{-12}$ (\code{ZERO\_TOL}) est arrondi à zéro. Les principaux
et taux nominaux sont inchangés (la dépréciation est réelle, pas nominale).

\subsection{Phase 6 : marché du crédit}
\label{sec:market}

\paragraph{Pool admissible et intensité $\eta$.}
Le pool admissible du pas $t$ est
\begin{equation}
N^{\mathrm{mkt}}_t=\#\{\text{vivantes n'ayant pas fait défaut au service
du pas courant}\},
\end{equation}
figé pendant toute la phase. Le nombre de tentatives d'appariement
(\emph{rounds}) est
\begin{equation}
R_t=
\begin{cases}
\bigl\lfloor \eta\bigl(N^{\mathrm{mkt}}_t\bigr)\bigr\rfloor
& \text{si } N^{\mathrm{mkt}}_t\ge2,\\[2pt]
0 & \text{sinon,}
\end{cases}
\qquad\boxed{\eta(N)=N}\ \text{(baseline M4.2)} .
\end{equation}
$\eta$ est implémentée comme une fonction nommée unique
(\code{model.eta}). Le cas $N^{\mathrm{mkt}}_t<2$ lève une ambiguïté
interne du cahier des charges (la formule $\lfloor\eta(N)\rfloor$
donnerait un round à $N=1$, mais la prose impose \og aucun round si le
pool contient moins de deux entités \fg) : le code suit la prose. M4B
exécutait à $N=1$ des rounds stériles via \code{permutation(1)}, qui ne
consomme aucun état RNG (vérifié à l'audit du 27 juillet 2026) : les deux
conventions sont donc \emph{équivalentes en dynamique et en flux
aléatoire} ; seule la valeur du compteur \code{mkt\_rounds} (absent de
M4B) diffère.

\paragraph{Formation d'une paire.}
Chaque round tire uniformément \emph{deux} identifiants distincts du pool,
sans remise à l'intérieur de la paire ($k\equiv2$). Le tirage préserve la
séquence RNG de M4B (\code{\_sample}) : permutation partielle si
$2\times8\ge N$, sinon tirages avec rejet. Les rounds sont indépendants et
une entité peut participer à plusieurs rounds. Si les deux capitaux
diffèrent, la plus riche devient prêteuse $\ell$, la plus pauvre
emprunteuse $b$ ; si $K_\ell\le K_b$ (égalité), le round est stérile.

\paragraph{Taux d'intérêt issu de la concavité.}
Les deux rendements marginaux sont évalués par \eqref{eq:marginal} (avec
plancher $K\mapsto\max(K,\code{K\_FLOOR})$, §\ref{sec:tol}) :
$m_\ell=A\gamma K_\ell^{\gamma-1}$, $m_b=A\gamma K_b^{\gamma-1}$, avec
$m_\ell<m_b$ puisque $K_\ell>K_b$ et $\gamma<1$. Le taux contractuel est la
\emph{moyenne géométrique}, institution de négociation explicite (et non
conséquence de la concavité) :
\begin{equation}
r_{\ell b}=\sqrt{m_\ell\,m_b}
= A\,\gamma\,(K_\ell K_b)^{(\gamma-1)/2}.
\label{eq:rate}
\end{equation}
Elle est symétrique, multiplicativement invariante, et égalise les
distances relatives : $r_{\ell b}/m_\ell=m_b/r_{\ell b}$ — c'est le milieu
de l'intervalle $[m_\ell,m_b]$ dans l'espace logarithmique.

\paragraph{Cible productive et principal.}
Le capital qui égalise rendement marginal et taux contractuel est
\begin{equation}
K^{*}(r_{\ell b})=\Bigl(\frac{A\gamma}{r_{\ell b}}\Bigr)^{1/(1-\gamma)} .
\label{eq:target}
\end{equation}

\begin{quote}
\textbf{Identité (cible géométrique).} Avec le taux \eqref{eq:rate},
$K^{*}(r_{\ell b})=\sqrt{K_\ell K_b}$ pour tout $0<\gamma<1$ et tout $A>0$.

\emph{Démonstration.} En substituant \eqref{eq:rate} dans
\eqref{eq:target} :
\[
K^{*}
=\Bigl(\frac{A\gamma}{A\gamma\,(K_\ell K_b)^{(\gamma-1)/2}}\Bigr)^{1/(1-\gamma)}
=\Bigl((K_\ell K_b)^{(1-\gamma)/2}\Bigr)^{1/(1-\gamma)}
=(K_\ell K_b)^{1/2}. \qquad\square
\]
\end{quote}

Conséquence à garder à l'esprit dans l'analyse : à capitaux donnés, la
cible instantanée $\sqrt{K_\ell K_b}$ ne dépend pas de $\gamma$ ; $\gamma$
agit sur la production, les taux \eqref{eq:rate} et, indirectement, sur les
capitaux et le réseau au cours du temps. L'identité est testée
numériquement à $10^{-12}$ relatif sur une grille de capitaux, de $\gamma$
et de $A$ (\code{tests/test\_engine.py}).

Le principal transféré est
\begin{equation}
q=\min\bigl(K_\ell-K^{*},\ \max(0,\,K^{*}-K_b)\bigr) ;
\end{equation}
si $q<\code{MIN\_LOAN}=10^{-9}$ le round est abandonné. Sinon
$K_\ell\leftarrow K_\ell-q$, $K_b\leftarrow K_b+q$, et le contrat
$(\ell,b,q,r_{\ell b})$ est créé ou fusionné. Comme $K_\ell>K^{*}>K_b$
(moyenne géométrique de deux valeurs distinctes), on a $q>0$ au sens réel ;
le transfert conserve le capital réel total et la valeur nette de chacune
des deux parties (testé).

\paragraph{Fusion des contrats.}
\label{sec:merge}
Si la paire orientée $(\ell,b)$ possède déjà un contrat $(q_0,r_0)$, le
nouveau prêt $(q_1,r_1)$ fusionne en
\begin{equation}
q\leftarrow q_0+q_1,\qquad
r\leftarrow \frac{q_0 r_0+q_1 r_1}{q_0+q_1},
\end{equation}
ce qui préserve exactement le flux d'intérêts total $r\,q=q_0r_0+q_1r_1$
(testé à $10^{-12}$). La paire inverse $(b,\ell)$ reste un contrat
distinct.

\paragraph{Compteurs de fonctionnement de $\eta$.}
La phase de marché enregistre séparément, par pas :
$N^{\mathrm{mkt}}_t$ (\code{mkt\_pool}), $R_t$ (\code{mkt\_rounds},
égal au nombre de rounds \emph{exécutés} — zéro si $N^{\mathrm{mkt}}_t<2$),
le nombre de transactions réussies (\code{new\_loans}, sémantique M4B), le
nombre de nouvelles arêtes contractuelles (\code{mkt\_new\_edges}), le
nombre de fusions (\code{mkt\_merges}) et le volume de principal échangé
(\code{loan\_volume}), avec l'invariant
$\code{new\_loans}=\code{mkt\_new\_edges}+\code{mkt\_merges}$ (testé).

\subsection{Phase 7 : faillites et cascades}

Règle M4B conservée sans modification de principe. La file initiale
(\emph{racines}) contient les vivantes en défaut de service du pas
(\code{liquidity}), de valeur nette $\NW_i<-\code{INSOLVENCY\_TOL}$
(\code{insolvency}), ou les deux (\code{both}), triées par identifiant.
Chaque itération (\emph{génération}) tue toutes les entités de la file :
\begin{itemize}
\item chaque prêteuse de la morte perd \emph{intégralement} son principal
      (perte de créance ; une arête de perte morte$\to$prêteuse est
      enregistrée) ;
\item les créances portées par la morte sont annulées (gain net pour ses
      emprunteuses, dont la dette disparaît) ;
\item le capital résiduel de la morte est détruit ;
\end{itemize}
puis la file suivante est reconstituée : vivantes avec
$\NW<-\code{INSOLVENCY\_TOL}$ (triées par identifiant). La résolution
itère jusqu'au point fixe (plus aucune insolvable), avec un garde-fou à
$n_{\mathrm{vivantes}}+1$ itérations (inatteignable : toute itération non
vide tue au moins une entité).

Conventions de résolution, héritées de M4B et explicitées à l'audit :
\begin{itemize}
\item une entité en file au début d'une génération meurt
\emph{inconditionnellement}, même si la mort d'une co-listée annule une
de ses dettes au cours de la même génération et l'aurait sauvée ; seul le
soulagement issu d'une génération \emph{antérieure} propage (le résultat
est indépendant de l'ordre des identifiants) ;
\item un défaut de service tue en génération 1 (cause \code{liquidity})
même si la valeur nette reste positive ;
\item la mort d'une \emph{prêteuse} n'enregistre aucune perte :
l'annulation de ses créances efface la dette de ses emprunteuses sans
flux (aucune partie vivante ne perd) ; \code{claim\_losses} ne compte que
le principal perdu par les prêteuses d'une \emph{emprunteuse} morte, et
$\sum_{\text{avalanches du pas}}V=\code{claim\_losses}$ du pas (vérifié à
$9{\cdot}10^{-12}$ près) ;
\item les nouveau-nées du pas subissent choc, production et dépréciation
dès leur pas de naissance.
\end{itemize}

\paragraph{Avalanches causales.}
Une avalanche est une composante faiblement connexe du graphe des arêtes de
perte \emph{restreint aux entités mortes pendant la même résolution}
(union--find). Par avalanche sont enregistrés : taille $s$ (nombre de
membres), profondeur (génération maximale de décès), nombre de racines,
causes, membres, et volume
\begin{equation}
V=\sum_{\text{arêtes de perte dont la source est membre}} q ,
\end{equation}
qui inclut les pertes infligées à des prêteuses survivantes (l'origine du
lien affecte la perte à la cascade qui l'a provoquée). La définition est
identique à M4B et ne change pas pendant l'étude.

\subsection{Phase 8 : mesures}

Aucun tirage aléatoire. La ligne de série du pas contient les 23 colonnes
M4B (démographie, capital, production, crédit, intérêts, racines par
cause, itérations de cascade, avalanches, pertes, injection, dépréciation,
gain de choc) plus les quatre compteurs de marché
(§\ref{sec:market}). Le statut passe à \code{extinction} si la population
est nulle, à \code{explosion} si elle dépasse \code{pop\_max} (ce qui
arrête le run mais ne constitue pas, en soi, un diagnostic d'explosion —
voir le protocole).

\section{Tolérances numériques}
\label{sec:tol}

Toutes héritées de M4B, à l'exception du plancher nommé :
\begin{center}
\begin{tabular}{lll}
\toprule
Constante & Valeur & Rôle\\
\midrule
\code{INSOLVENCY\_TOL} & $10^{-9}$ & seuil d'insolvabilité sur $\NW$\\
\code{MIN\_LOAN} & $10^{-9}$ & principal minimal d'un prêt\\
\code{ZERO\_TOL} & $10^{-12}$ & arrondi à zéro du capital après dépréciation\\
\code{K\_FLOOR} & $10^{-9}$ & plancher de $K$ dans $m_\gamma(K)$\\
\bottomrule
\end{tabular}
\end{center}

\code{K\_FLOOR} est la même valeur que le $\max(K,10^{-9})$ écrit en dur
dans \code{\_pair\_terms} de M4B ; M4.2 le nomme parce que son poids croît
quand $\gamma$ décroît : au plancher, le rendement marginal vaut
$A\gamma\,10^{9(1-\gamma)}$ (environ $3{\cdot}10^{5}$ à $\gamma=1/3$ contre
$1{,}6{\cdot}10^{4}$ à $\gamma=1/2$). La cible $K^*=\sqrt{K_\ell K_b}$
(avec capitaux planchonnés) reste bornée par les capitaux réels et $q$ est
plafonné par $K_\ell-K^{*}$ : aucun débordement flottant n'est possible en
double précision. Le comportement au plancher est testé
(\code{test\_k\_floor\_is\_applied\_and\_finite}).

$\gamma=1$ est une singularité de \eqref{eq:target} ($1/(1-\gamma)$) et de
\eqref{eq:fixedpoint} : la validation de la configuration impose
$\gamma<1$ strictement, sans garde supplémentaire dans le code.

\paragraph{Fuites d'arrondi héritées, mesurées à l'audit du 27 juillet
2026.} Deux imprécisions structurelles, identiques dans M4B et conservées
pour la parité, sont documentées comme tolérances :
\begin{itemize}
\item \emph{puits de dépréciation} : quand $(1-\delta)K\le\code{ZERO\_TOL}$,
le capital est planché à zéro et le reliquat ($\le10^{-12}$ par entité
planchée) n'est pas ajouté à la dépréciation mesurée ; zéro déclenchement
observé en régime normal (run d'audit $\gamma=0{,}6$, $T=300$) ;
\item \emph{dérive du service agrégé} : le débit utilise le cache
incrémental $S_b$ tandis que les crédits recalculent $\sum q\,r$ ; écart
mesuré $\le5{,}5\cdot10^{-13}$ par pas, cumul $\sim10^{-11}$ sur 300 pas.
\end{itemize}
Résidu de bilan mesuré : $1{,}3\cdot10^{-11}$ max par pas (seuil de test
$10^{-6}$). Les dictionnaires d'agrégats conservent une clé par entité
née (jamais purgée, héritage M4B) : coût mémoire $O(\text{naissances})$,
sans effet sur les valeurs (résidus $\le1{,}5\cdot10^{-12}$, contrôlés par
\code{consistency\_errors} au seuil $10^{-6}$).

\section{Point fixe autarcique discret et analyse dimensionnelle}
\label{sec:fixedpoint}

Sans choc, sans crédit et sans naissance, la mise à jour exécutée d'un pas
est $K\mapsto(1-\delta)\bigl(K+A K^\gamma\bigr)$ (la production
\emph{précède} la dépréciation). Le point fixe exact de cette relation
\emph{discrète} — à ne pas remplacer par son approximation continue —
satisfait $K=(1-\delta)(K+AK^\gamma)$, d'où
\begin{equation}
\Kaut(\gamma)=\Bigl(\frac{(1-\delta)\,A}{\delta}\Bigr)^{1/(1-\gamma)} .
\label{eq:fixedpoint}
\end{equation}
À la configuration centrale ($\delta=0{,}05$, $A=1$), la base vaut $19$ :
\begin{center}
\begin{tabular}{lccccc}
\toprule
$\gamma$ & $1/3$ & $0{,}4$ & $1/2$ & $0{,}6$ & $2/3$\\
$\Kaut$ (J) & $82{,}8$ & $135{,}3$ & $361$ & $1573{,}6$ & $6859$\\
\bottomrule
\end{tabular}
\end{center}
La convergence locale est géométrique de raison
$1-\delta(1-\gamma)$ : lente à grand $\gamma$ (raison $0{,}983$ à
$\gamma=2/3$), ce qui gouverne les transitoires ($K_0=25\ll\Kaut$). Le
moteur atteint \eqref{eq:fixedpoint} à $10^{-9}$ relatif pour
$\gamma\in\{1/3,1/2,2/3\}$ (\code{test\_autarkic\_fixed\_point}).

Deux identités dimensionnelles guident les prédictions de campagne :
\begin{equation}
\frac{F_\gamma(\Kaut)}{\Kaut}=\frac{\delta}{1-\delta}
\quad\text{(invariante en $\gamma$)},\qquad
m_\gamma(\Kaut)=\gamma\,\frac{\delta}{1-\delta}
\quad\text{(linéaire en $\gamma$)} .
\label{eq:dimensional}
\end{equation}
À l'échelle stationnaire, le taux d'intérêt typique croît donc comme
$\gamma$ tandis que la production par unité de capital est fixe : la charge
d'intérêts d'une emprunteuse au levier $D/K$, rapportée à sa production,
vaut environ $\gamma\,(D/K)$. Prédictions testables (exploratoires) :
$\gamma$ élevé $\Rightarrow$ service relativement plus lourd, davantage de
racines de liquidité, vies plus courtes ; $\gamma$ faible $\Rightarrow$
crédit bon marché, moins d'avalanches (risque d'identifiabilité) ; l'échelle
$\Kaut$ varie de deux décades sur la grille pilote, d'où le contrôle
d'échelle apparié ci-dessous.

\section{Contrôle d'échelle apparié $A_{\mathrm{norm}}$}
\label{sec:anorm}

Faire varier $\gamma$ à $A=1$ change simultanément la courbure et l'échelle
\eqref{eq:fixedpoint}. L'ablation pré-enregistrée fixe un capital de
référence $K_{\mathrm{ref}}=\Kaut(1/2)=361$ J et utilise, uniquement dans
ce bras,
\begin{equation}
A_{\mathrm{norm}}(\gamma)=\frac{\delta}{1-\delta}\,
K_{\mathrm{ref}}^{\,1-\gamma}=19^{\,1-2\gamma}
\qquad(\delta=0{,}05),
\end{equation}
qui impose $\Kaut(\gamma; A_{\mathrm{norm}})=K_{\mathrm{ref}}$ pour tout
$\gamma$ (sanité : $A_{\mathrm{norm}}(1/2)=1$ ; $A_{\mathrm{norm}}(1/3)
= 19^{1/3}\approx2{,}668$ ; $A_{\mathrm{norm}}(2/3)=19^{-1/3}\approx
0{,}375$). $A$ est un champ de configuration tracé précisément pour que ce
bras ne soit pas un second modèle silencieux ; le protocole fige $A=1$
partout ailleurs.

\paragraph{Invariance d'échelle exacte (propriété structurelle, vérifiée
empiriquement le 27 juillet 2026).} À $(\gamma,\delta,\sigma,\lambda)$
fixés, la transformation
\begin{equation}
K_i \mapsto c\,K_i\ \ (\forall i),\qquad K_0\mapsto c\,K_0,\qquad
A\mapsto c^{\,1-\gamma}A,\qquad c>0,
\end{equation}
laisse \emph{exactement} invariante toute l'histoire discrète du modèle
pour une même suite de tirages aléatoires : naissances, chocs
$\xi_{i,t}$ (multiplicatifs, insensibles à $c$), comparaisons
prêteuse/emprunteuse (préservées par toute mise à l'échelle positive),
défauts de service (le ratio $\varphi_b=K_b/S_b$ est invariant car $K_b$
et le service dû $S_b$ scalent tous deux par $c$ --- le taux
$r_{\ell b}=A\gamma(K_\ell K_b)^{(\gamma-1)/2}$ est lui-même
\emph{inchangé} sous cette transformation, comme il se doit pour une
grandeur de dimension $1/\text{temps}$), insolvabilité, composition des
cascades. Seules les grandeurs \emph{extensives} (capitaux, principaux,
volumes) sont multipliées par $c$ ; toutes les grandeurs \emph{discrètes
ou sans dimension} (identités des mortes, tailles et profondeurs
d'avalanches, $\hat\tau$, $b$, $N/\lambda$) sont rigoureusement
identiques. Vérifié numériquement sur $\gamma\in\{1/3,1/2,0{,}6,2/3\}$
et $c\in\{0{,}1;4;7;19\}$ : écart relatif maximal
$1{,}2\cdot10^{-11}$, entièrement attribuable aux tolérances numériques
absolues (\code{K\_FLOOR}, \code{ZERO\_TOL}, \code{INSOLVENCY\_TOL}),
dix ordres de grandeur sous les capitaux en jeu dans la plage testée.

Puisque $\Kaut(\gamma,A)$ scale lui-même exactement par $c$ sous cette
transformation (immédiat depuis \eqref{eq:fixedpoint} avec
$A\to c^{1-\gamma}A$), le rapport adimensionnel $K_0/\Kaut(\gamma,A)$
est \emph{invariant}. Conséquence directe : à
$(\gamma,\delta,\sigma,\lambda)$ fixés, toute statistique sans dimension
du modèle (en particulier $\hat\tau$) ne peut dépendre de $A$ et $K_0$
\emph{qu'à travers} ce rapport --- ce n'est pas seulement un motif
empirique observé dans l'ablation, mais une conséquence nécessaire de
cette symétrie exacte de la dynamique. L'ablation à $A_{\mathrm{norm}}$
mesure donc, par construction, l'effet résiduel \emph{pur} de la
courbure $\gamma$ à rapport $K_0/\Kaut$ constant.

\section{Invariants et contrôles}

\begin{itemize}
\item \textbf{Bilan réel du pas} : en notant $I$ l'injection, $G$ le gain
de choc, $\Pi$ la production totale, $\Delta$ la dépréciation et $X$ le
capital détruit aux faillites,
$\;\Delta K_{\mathrm{tot}}=I+G+\Pi-\Delta-X$, vérifié à $10^{-6}$ relatif
par pas (les intérêts et transferts de principal sont des flux internes de
somme nulle). Testé à $\gamma=0{,}4$ et $0{,}65$.
\item \textbf{Créances $=$ dettes} : chaque contrat compte une fois de
chaque côté ; $\sum_i C_i=\sum_i D_i$ structurellement.
\item \textbf{Cohérence du carnet} : \code{consistency\_errors} reconstruit
$C$, $D$, $S$ depuis les contrats (tolérance $10^{-6}$ relative), vérifie
qu'aucun contrat n'implique une morte ni n'est dégénéré ; appelé hors
boucle (jamais dans \code{step()}), notamment dans le résumé de run.
\item \textbf{Reproductibilité} : même configuration $\Rightarrow$ mêmes
octets dans les tables (testé) ; les options d'enregistrement
(\code{snapshot\_every}, \code{individual\_every}) ne modifient pas les
trajectoires (testé).
\end{itemize}

\section{Parité mesurée avec M4B à $k=2$}
\label{sec:parity}

Le cas $\gamma=1/2$, $A=1$ doit retrouver M4B configuré à $k=2$, mêmes
paramètres et mêmes graines. Les formes génériques
$A\gamma K^{\gamma-1}$ et $(A\gamma/r)^{1/(1-\gamma)}$ peuvent différer
d'un ULP des formes M4B $1/(2\sqrt K)$ et $(1/(2r))^2$ ; conformément au
cahier des charges, aucune branche de compatibilité ne fabrique l'égalité
bit à bit : l'écart est mesuré et borné.

\textbf{Faits observés} (\code{tests/test\_parity\_m4b.py}, exécuté le
27 juillet 2026, moteur \code{m4\_2-1}) :
\begin{itemize}
\item régimes de \emph{gating} (graine 0, $\lambda=10$, $\sigma=0{,}25$,
120 pas ; graine 7, $\lambda=18$, $\sigma=0{,}31$, 90 pas) : tout le
discret est identique pas à pas (naissances, morts et leurs causes,
contrats et paires orientées, défauts, itérations de cascade, tailles,
profondeurs, racines et membres des avalanches, compteurs entiers des
séries) ; écart flottant relatif maximal $5{,}7\cdot10^{-13}$ et
$3{,}0\cdot10^{-13}$ (tolérance $10^{-9}$) sur capitaux, agrégats du
carnet, contrats et colonnes flottantes ;
\item sonde longue non bloquante (graine 1, $\lambda=30$, $T=1000$) :
\emph{aucune divergence discrète en 1000 pas} ; écart flottant relatif
maximal des séries $3{,}1\cdot10^{-12}$.
\end{itemize}
\textbf{Inférence} : les écarts d'arrondi entre les deux écritures ne
s'amplifient pas en bifurcation discrète sur ces horizons ; la parité est
établie au sens du cahier des charges (égalité logique exacte, égalité
numérique bornée par $\sim10^{-12}$ sur $10^3$ pas). L'unique écart de
règle (aucun round de marché si le pool $<2$, §\ref{sec:market}) n'a pas
été exercé dans ces runs (population toujours $\ge2$ au marché).

\section{Données enregistrées}

Par run : \code{config.json} (version du moteur, paramètres — dont
$\gamma$ et $A$ —, règles fixes, options d'enregistrement) ;
\code{series.csv} (27 colonnes : les 23 de M4B $+$ \code{mkt\_pool},
\code{mkt\_rounds}, \code{mkt\_new\_edges}, \code{mkt\_merges}) ;
\code{avalanches.csv} (identifiant, pas, taille, profondeur, racines,
volume en joules, causes) ; \code{avalanche\_members.csv} (appartenance,
racine, génération) ; \code{deaths.csv} (état terminal complet de chaque
morte) ; \code{entities.csv} (état civil) ; \code{final\_loans.csv}
(carnet final) ; \code{individual\_series.csv.gz} (trajectoires
longitudinales, fréquence réglable) ; \code{snapshots/entities\_t*.npz} et
\code{network\_t*.npz} (instantanés) ; \code{summary.json} (statut,
comptages, rapport de branchement, erreurs de carnet). Le schéma M4B est
conservé à l'identique (colonnes ajoutées en fin uniquement), condition de
réutilisation des recettes de figures Simulation Lab.

\section{Traçabilité}

Moteur : \code{m4\_2\_credit\_soc/m4\_2/model.py} et \code{io.py}, version
\code{m4\_2-1}, copie adaptée ligne à ligne de
\code{m4b\_credit\_soc\_mini/m4b/} (version \code{m4b-mini-3}, lecture
seule). Tests : \code{tests/test\_engine.py} (18 tests, points 1--13 du
cahier des charges), \code{tests/test\_parity\_m4b.py} (parité et sonde),
\code{tests/test\_stats.py} (point 14, estimateurs sur données
synthétiques). Les condensats du moteur sont inscrits dans les manifestes
de campagne. En cas de contradiction entre ce document et le code, le code
exécuté fait foi et l'écart doit être documenté.

\end{document}
